


% \documentclass[preprint,12pt,authoryear]{elsarticle}
%% Use the option review to obtain double line spacing
%% \documentclass[authoryear,preprint,review,12pt]{elsarticle}

%% Use the options 1p,twocolumn; 3p; 3p,twocolumn; 5p; or 5p,twocolumn
%% Use the options 1p,twocolumn; 3p; 3p,twocolumn; 5p; or 5p,twocolumn
%% for a journal layout:
%% \documentclass[final,1p,times,authoryear]{elsarticle}
% \documentclass[final,1p,times,twocolumn,authoryear]{elsarticle}
% \documentclass[final,3p,times,authoryear]{elsarticle}
%% \documentclass[final,3p,times,twocolumn,authoryear]{elsarticle}
%% \documentclass[final,5p,times,authoryear]{elsarticle}
\documentclass[final,5p,times,twocolumn,authoryear]{elsarticle}
%% For including figures, graphicx.sty has been loaded in
%% elsarticle.cls. If you prefer to use the old commands
%% please give \usepackage{epsfig}

%% The amssymb package provides various useful mathematical symbols
\usepackage{amssymb}
%% The amsmath package provides various useful equation environments.
\usepackage{amsmath}
%% The amsthm package provides extended theorem environments
%% \usepackage{amsthm}

%% The lineno packages adds line numbers. 
\usepackage{lineno}
% Start line numbering with
%% \begin{linenumbers}, end it with \end{linenumbers}. Or switch it on
%% for the whole article with \linenumbers.


\usepackage{tabularx}
\usepackage{kotex}
\usepackage{booktabs}
\usepackage{supertabular}
\usepackage{afterpage}
\usepackage{graphicx}
\graphicspath{{./images/}}
% hyperlink related 
\usepackage{url}       % URL 줄바꿈을 더 깔끔하게 처리
\usepackage{hyperref}  % 클릭 가능한 링크 생성


% 캡션 정렬.
\usepackage{caption}
\captionsetup[figure]{labelsep=period}

\journal{Applied Energy}

\begin{document}

\begin{frontmatter}




%% Title, authors and addresses

%% use the tnoteref command within \title for footnotes;
%% use the tnotetext command for theassociated footnote;
%% use the fnref command within \author or \affiliation for footnotes;
%% use the fntext command for theassociated footnote;
%% use the corref command within \author for corresponding author footnotes;
%% use the cortext command for theassociated footnote;
%% use the ead command for the email address,
%% and the form \ead[url] for the home page:

%% \title{Title\tnoteref{label1}}
%% \tnotetext[label1]{}
%% \author{Name\corref{cor1}\fnref{label2}}
%% \ead{email address}
%% \ead[url]{home page}
%% \fntext[label2]{}
%% \cortext[cor1]{}
%% \affiliation{organization={},
%%            addressline={}, 
%%            city={},
%%            postcode={}, 
%%            state={},
%%            country={}}
%% \fntext[label3]{}


\title{Contextually Extended iTransformer (CEiT) for Query-Based Virtual Sensing in Dynamic Spatial Sensor Networks} %% Article title

%% 추가 후보 
% Contextually Extended iTransformer for Robust Virtual Sensing with Missing and Variable Sensors

% Contextually Extended iTransformer for Virtual Sensing of Indoor Air Quality in Livestock Buildings



%% 저자 및 소속 설정
\author[a]{Habin Jo}

\author[b]{Masanori Shukuya}

\author[a,c]{Wonjun Choi}
\ead{wonjun.choi@jnu.ac.kr} %% 마지막 저자의 이메일

%% 하단 주석 설정
\cortext[cor1]{Corresponding author}

%% 소속 정보 (이전과 동일)
\affiliation[a]{
    organization={Department of Architecture and Civil Engineer},
    addressline={77 Yongbong-ro, Buk-gu},
    city={Gwangju},
    postcode={61186},
    country={South Korea}
}

\affiliation[b]{organization={Department of Restoration Ecology and Built Environment, Tokyo City University},
            addressline={Address B},
            city={Yokohama},
            postcode={67890},
            country={Japan}}

\affiliation[c]{organization={Department of Architecture and Civil Engineering, Chonnam National University}, 
            addressline={77 Yongbong-ro, Buk-gu}, 
            city={Gwangju}, 
            postcode={61186},
            country={South Korea}}

\begin{abstract}
%% Text of abstract
abstract 
for real 







\end{abstract}

%%Graphical abstract
% \begin{graphicalabstract}
%\includegraphics{grabs}
% \end{graphicalabstract}

%%Research highlights
% \begin{highlights}
% \item Research highlight 1
% \item Research highlight 2
% \end{highlights}

%% Keywords
\begin{keyword}

Virtual sensor \sep Deep learning \sep Transformer \sep Smart farm \sep  Spatio-temporal interpolation \sep Interpretable AI 

%% keywords here, in the form: keyword \sep keyword

%% PACS codes here, in the form: \PACS code \sep code

%% MSC codes here, in the form: \MSC code \sep code
%% or \MSC[2008] code \sep code (2000 is the default)

\end{keyword}

\end{frontmatter}



%% Add \usepackage{lineno} before \begin{document} and uncomment 
%% following line to enable line numbers
%% \linenumbers




\renewcommand{\figurename}{Fig.} % 메인 텍스트용 figure 표시 설정

%%%%%%%%%%%%%%%%%%%%% MAIN TEXT 
%% Use \section commands to start a section
% \begin{linenumbers}
% \modulolinenumbers[2] % 5, 10, 15... 행에만 번호 표시

\section{Introduction}
\label{sec:intro}

%% Labels are used to cross-reference an item using \ref command.

\subsection{Building Decarbonization and Electrification Trends}
전 세계적으로 건물 부문은 에너지 사용의 약 32\%, CO$_2$ 배출의 34\%를 차지하며, 기후 변화 대응을 위한 핵심 감축 대상으로 지목되고 있음 \cite{unep2024global}. 이에 따라 화석 연료 기반의 난방 및 급탕 시스템을 전기에너지 기반으로 전환하는 `건물 전전화(Building Electrification)'가 주요 정책 및 기술적 솔루션으로 부상하고 있음 \cite{irena2020renewable}.

\subsection{Importance of DHW Energy and the Rise of Heat Pumps}
건물의 단열 성능이 강화됨에 따라 냉난방 부하가 감소하는 반면, 급탕(DHW) 수요는 일정하게 유지되어 전체 건물 에너지 사용에서 차지하는 비중이 15--40\%까지 증가하고 있음 \cite{ref3, ref4}. 이러한 급탕 부문의 효율화를 위해 전통적인 가스 보일러(GB)를 고효율 히트펌프 보일러(HPB)로 대체하는 것이 가장 유력한 대안으로 여겨짐 \cite{ref3, ref5}.

\subsection{Problem Statement: Regional Variations in Heat Pump Effectiveness due to Grid Mix}
히트펌프는 사용 단계에서 높은 효율을 보이지만, 히트펌프의 투입되는 전력을 공급하는 그리드의 전력 발전 믹스에 따라, 전력의 CO$_2$ 배출 계수와  Primary energy factor가 크게 달라짐. 이는 히트펌프 보일러의 도입을 통해 탄소중립을 효과적으로 달성하기 위해서는 히트펌프 보일러의 기기의 효율(COP)를 높이는 것 뿐만 아니라, 재생에너지에 기반한 전력 발전의 비중을 높여 전력 사용에 의한 Primary energy use와 CO$_2$ 배출 계수를 낮추는 것이 선행되어함을 나타낸다.\cite{irena2020renewable}.

\subsection{Research Objective}
본 연구는 가스 보일러와 히트펌프 보일러의 성능을 한국의 발전믹스의 재생에너지 비중이 증가함에 따라 가스 보일러와 히트펌프 보일러의 primary energy use와 CO$_2$ 배출량을 비교하여, 한국 전력 그리드의 발전 비중 증가에 따른 히트펌프 보일러로의 전환이 에너지 및 환경에 미치는 영향을 평가한다. 


\section{System Setting}
\label{sec:system}

본 연구에서는 가스 보일러(GB)와 공기원 히트펌프 보일러(HPB) 두 가지 급탕 시스템을 동일한 급탕 사용 스케줄과 외기 온도 조건에서 비교 분석하였다.  Figure~\ref{fig:hot water use schedule}에 본 연구에서 사용한 최종 사용 온수의 공급 스케줄을 나타냈으며, 아침과 저녁 시간에 주된 급탕 사용이 발생하는 주거용 온수 사용 스케줄의 특징을 반영하였으며, 약 200 L의 사용량을 가정하였다. 

\begin{figure}[htbp]
  \centering
  \includegraphics[width = 9 cm]{images/water_use_schedule_ashpb.png}
  \captionsetup{width=\textwidth}
  \caption{본 연구에서 설정한 급탕 사용 스케줄}
  \label{fig:hot water use schedule}
\end{figure}

Figures~\ref{fig:gb_diagram}--\ref{fig:hpb_diagram}은 각각 가스 보일러와 히트펌프 보일러의 시스템의 구조를 나타내며, energy flow rate, 온도, 체적유량 등의 symbol들이 표기되어있다. 가스 보일러(Figure~\ref{fig:gb_diagram})는, Combustion chamber로 유입된 상수도를 60 °C로 가열해 mixing valve로 공급하며, mixing valve에서 15 °C의 상수도와 혼합되어 40 °C의 온수를 온수 사용 스케줄(Figure~\ref{fig:hot water use schedule})에 따라 공급한다.
이때 혼합비($\alpha$)는 다음과 같이 계산된다:

\begin{equation}
\label{eq:gb_mixing_ratio}
\alpha = \frac{T_{\text{serv,w}} - T_{\text{sup,w}}}{T_{\text{comb}} - T_{\text{sup,w}}}
\end{equation}

여기서 $\alpha$는 보일러 출수 비율이며, $(1-\alpha)$는 상수도 직접 공급 비율이다. 유량 관계는 $\dot{V}_{w,\text{sup,comb}} = \alpha \dot{V}_{w,\text{serv}}$ 및 $\dot{V}_{w,\text{sup,mix}} = (1 - \alpha) \dot{V}_{w,\text{serv}}$로 표현된다.

\begin{figure*}[htbp]
  \centering
  \includegraphics[width=10 cm]{images/GB_diagram.png}
  \captionsetup{width=\textwidth}
  \caption{Gas boiler system diagram.}
  \label{fig:gb_diagram}
\end{figure*}

\begin{figure*}[htbp]
  \centering
  \includegraphics[width = 17 cm]{images/HPB_diagram.png}
  \captionsetup{width=\textwidth}
  \caption{Air source heat pump boiler system diagram.}
  \label{fig:hpb_diagram}
\end{figure*}
\subsection{Gas Boiler System Modeling}

가스 보일러 시스템은 연소실(combustion chamber)과 믹싱밸브(mixing valve) 두 개의 주요 서브시스템으로 구성된다. 각 서브시스템에 대해 에너지, 엔트로피, 엑서지 분석을 수행한다.

\subsubsection{Combustion Chamber}

연소실에서는 천연가스($E_{\text{NG}}$)를 연소하여 급수($\dot{V}_{w,\text{sup,comb}}$)를 설정 온도($T_{\text{comb}}$)로 가열한다. 연소 효율($\eta_{\text{comb}}$)을 고려한 에너지 밸런스는 다음과 같이 표현된다:

\begin{equation}
\label{eq:gb_energy_comb}
E_{\text{NG}} = \frac{Q_{\text{comb,load}}}{\eta_{\text{comb}}}
\end{equation}

\begin{equation}
\label{eq:gb_energy_balance_comb}
E_{\text{NG}} + Q_{w,\text{sup}} = Q_{w,\text{comb}} + Q_{\text{exh}}
\end{equation}

여기서 $Q_{\text{comb,load}}$는 보일러 유량 기준 필요한 열량[W]이며, 다음과 같이 계산된다:


\begin{equation}
\label{eq:gb_comb_load}
Q_{\text{comb,load}} = c_w \rho_w \dot{V}_{w,\text{sup,comb}} (T_{\text{comb}} - T_{\text{sup,w}})
\end{equation}

급수 열량($Q_{w,\text{sup}}$), 보일러 출수 열량($Q_{w,\text{comb}}$), 배기가스 열량($Q_{\text{exh}}$)은 각각 다음과 같이 계산된다:

\begin{equation}
\label{eq:gb_qw_sup}
Q_{w,\text{sup}} = c_w \rho_w \dot{V}_{w,\text{sup,comb}} (T_{\text{sup,w}} - T_0)
\end{equation}

\begin{equation}
\label{eq:gb_qw_comb}
Q_{w,\text{comb}} = c_w \rho_w \dot{V}_{w,\text{sup,comb}} (T_{\text{comb}} - T_0)
\end{equation}

\begin{equation}
\label{eq:gb_qexh}
Q_{\text{exh}} = (1 - \eta_{\text{comb}}) E_{\text{NG}}
\end{equation}

연소실에서의 엔트로피 생성은 다음과 같이 계산된다. 천연가스의 엔트로피는:

\begin{equation}
\label{eq:gb_s_ng}
S_{\text{NG}} = \frac{E_{\text{NG}}}{T_{\text{NG}}}
\end{equation}

여기서 $T_{\text{NG}} = T_0 / (1 - \eta_{\text{X,NG}})$이며, $\eta_{\text{X,NG}} = 0.93$은 천연가스의 엑서지 효율이다. 급수, 보일러 출수, 배기가스의 엔트로피는 각각:

\begin{equation}
\label{eq:gb_s_w_sup}
S_{w,\text{sup}} = c_w \rho_w \dot{V}_{w,\text{sup,comb}} \ln\left(\frac{T_{\text{sup,w}}}{T_0}\right)
\end{equation}

\begin{equation}
\label{eq:gb_s_w_comb}
S_{w,\text{comb}} = c_w \rho_w \dot{V}_{w,\text{sup,comb}} \ln\left(\frac{T_{\text{comb}}}{T_0}\right)
\end{equation}

\begin{equation}
\label{eq:gb_s_exh}
S_{\text{exh}} = \frac{Q_{\text{exh}}}{T_{\text{exh}}}
\end{equation}

연소실에서의 엔트로피 생성($S_{g,\text{comb}}$)은 다음과 같이 계산된다:

\begin{equation}
\label{eq:gb_sg_comb}
S_{g,\text{comb}} = (S_{w,\text{comb}} + S_{\text{exh}}) - (S_{\text{NG}} + S_{w,\text{sup}})
\end{equation}

연소실에서의 엑서지 분석을 수행하면, 천연가스 엑서지 입력은:

\begin{equation}
\label{eq:gb_x_ng}
X_{\text{NG}} = \eta_{\text{X,NG}} E_{\text{NG}}
\end{equation}

급수, 보일러 출수, 배기가스의 엑서지는 각각:

\begin{equation}
\label{eq:gb_x_w_sup}
X_{w,\text{sup}} = c_w \rho_w \dot{V}_{w,\text{sup,comb}} \left[(T_{\text{sup,w}} - T_0) - T_0 \ln\left(\frac{T_{\text{sup,w}}}{T_0}\right)\right]
\end{equation}

\begin{equation}
\label{eq:gb_x_w_comb}
X_{w,\text{comb}} = c_w \rho_w \dot{V}_{w,\text{sup,comb}} \left[(T_{\text{comb}} - T_0) - T_0 \ln\left(\frac{T_{\text{comb}}}{T_0}\right)\right]
\end{equation}

\begin{equation}
\label{eq:gb_x_exh}
X_{\text{exh}} = \left(1 - \frac{T_0}{T_{\text{exh}}}\right) Q_{\text{exh}}
\end{equation}

연소실에서의 엑서지 소비($X_{c,\text{comb}}$)는 엔트로피 생성과 관련하여 다음과 같이 계산된다:

\begin{equation}
\label{eq:gb_xc_comb}
X_{c,\text{comb}} = T_0 S_{g,\text{comb}}
\end{equation}

\subsubsection{Mixing Valve}

믹싱밸브에서는 보일러 출수($T_{\text{comb}}$)와 상수도($T_{\text{sup,w}}$)를 혼합하여 서비스 급탕 온도($T_{\text{serv,w}}$)로 공급한다. 에너지 밸런스는 다음과 같이 표현된다:

\begin{equation}
\label{eq:gb_energy_mix}
Q_{w,\text{comb}} + Q_{w,\text{sup,mix}} = Q_{w,\text{serv}}
\end{equation}

서비스 급탕 온도는 혼합비를 이용하여 다음과 같이 계산된다:

\begin{equation}
\label{eq:gb_t_serv}
T_{\text{serv,w}} = \alpha T_{\text{comb}} + (1 - \alpha) T_{\text{sup,w}}
\end{equation}

믹싱밸브로 직접 공급되는 상수도 열량과 서비스 급탕 열량은 각각:

\begin{equation}
\label{eq:gb_qw_sup_mix}
Q_{w,\text{sup,mix}} = c_w \rho_w \dot{V}_{w,\text{sup,mix}} (T_{\text{sup,w}} - T_0)
\end{equation}

\begin{equation}
\label{eq:gb_qw_serv}
Q_{w,\text{serv}} = c_w \rho_w \dot{V}_{w,\text{serv}} (T_{\text{serv,w}} - T_0)
\end{equation}

믹싱밸브에서의 엔트로피 생성은 다음과 같이 계산된다:

\begin{equation}
\label{eq:gb_s_w_sup_mix}
S_{w,\text{sup,mix}} = c_w \rho_w \dot{V}_{w,\text{sup,mix}} \ln\left(\frac{T_{\text{sup,w}}}{T_0}\right)
\end{equation}

\begin{equation}
\label{eq:gb_s_w_serv}
S_{w,\text{serv}} = c_w \rho_w \dot{V}_{w,\text{serv}} \ln\left(\frac{T_{\text{serv,w}}}{T_0}\right)
\end{equation}

\begin{equation}
\label{eq:gb_sg_mix}
S_{g,\text{mix}} = S_{w,\text{serv}} - (S_{w,\text{comb}} + S_{w,\text{sup,mix}})
\end{equation}

믹싱밸브에서의 엑서지 분석을 수행하면, 서비스 급탕 엑서지는:

\begin{equation}
\label{eq:gb_x_w_serv}
X_{w,\text{serv}} = c_w \rho_w \dot{V}_{w,\text{serv}} \left[(T_{\text{serv,w}} - T_0) - T_0 \ln\left(\frac{T_{\text{serv,w}}}{T_0}\right)\right]
\end{equation}

믹싱밸브로 직접 공급되는 상수도 엑서지는:

\begin{equation}
\label{eq:gb_x_w_sup_mix}
X_{w,\text{sup,mix}} = c_w \rho_w \dot{V}_{w,\text{sup,mix}} \left[(T_{\text{sup,w}} - T_0) - T_0 \ln\left(\frac{T_{\text{sup,w}}}{T_0}\right)\right]
\end{equation}

믹싱밸브에서의 엑서지 소비($X_{c,\text{mix}}$)는:

\begin{equation}
\label{eq:gb_xc_mix}
X_{c,\text{mix}} = T_0 S_{g,\text{mix}}
\end{equation}

\subsubsection{System Exergy Efficiency}

가스 보일러 시스템 전체의 엑서지 효율($\eta_{\text{X,GB}}$)은 다음과 같이 정의된다:

\begin{equation}
\label{eq:gb_exergy_efficiency}
\eta_{\text{X,GB}} = \frac{X_{w,\text{serv}}}{X_{\text{NG}}}
\end{equation}

총 엑서지 소비는 연소실과 믹싱밸브에서의 엑서지 소비의 합으로 계산된다:

\begin{equation}
\label{eq:gb_xc_tot}
X_{c,\text{tot}} = X_{c,\text{comb}} + X_{c,\text{mix}}
\end{equation}

\subsection{Air Source Heat Pump Boiler System Modeling}

공기원 히트펌프 보일러 시스템은 실외기(outdoor unit), 냉매루프(refrigerant loop), 저탕조(hot water tank), 믹싱밸브로 구성된다. 본 연구에서는 냉매 상태를 최적화 알고리즘을 통해 결정하는 방법을 채택하였다.

\subsubsection{Optimization Problem}

히트펌프 사이클의 최적 운전점은 총 전력 사용률을 최소화하는 최적화 문제를 통해 결정된다:

\begin{equation}
\label{eq:hpb_optimization}
\begin{aligned}
\min_{\Delta T_{\text{ref,cond}}, \Delta T_{\text{ref,evap}}} \quad & E_{\text{tot}} = E_{\text{cmp}} + E_{\text{fan,ou}} \\
\text{subject to} \quad & Q_{\text{cond,load}} \leq Q_{\text{LMTD,cond}} \leq Q_{\text{cond,load}} \cdot (1 + \epsilon) \\
& Q_{\text{ref,evap}} \cdot (1 - \epsilon) \leq Q_{\text{LMTD,evap}} \leq Q_{\text{ref,evap}} \cdot (1 + \epsilon)
\end{aligned}
\end{equation}

여기서 $E_{\text{tot}}$는 총 전력 사용률[W], $E_{\text{cmp}}$는 압축기 전력[W], $E_{\text{fan,ou}}$는 실외기 팬 전력[W]이다. 최적화 변수는 $\Delta T_{\text{ref,cond}}$ (냉매-저탕조 온도차[K])와 $\Delta T_{\text{ref,evap}}$ (냉매-실외 공기 온도차[K])이며, $\epsilon$은 허용 오차(본 연구에서는 1\%)이다.

\subsubsection{Refrigerant Cycle States}

증발 온도와 응축 온도는 최적화 변수로부터 다음과 같이 계산된다:

\begin{equation}
\label{eq:hpb_temperatures}
T_{\text{evap}} = T_0 - \Delta T_{\text{ref,evap}}, \quad T_{\text{cond}} = T_{\text{tank,w}} + \Delta T_{\text{ref,cond}}
\end{equation}

냉매 사이클의 각 상태점(State 1-4)의 물성치는 증발 온도와 응축 온도를 기반으로 계산된다. State 1은 증발기 출구(포화 증기), State 2는 압축기 출구, State 3은 응축기 출구(포화 액체), State 4는 팽창밸브 출구이다.

냉매 유량($\dot{m}_{\text{ref}}$)은 목표 열 교환율을 만족하도록 계산된다:

\begin{equation}
\label{eq:hpb_refrigerant_flow}
\dot{m}_{\text{ref}} = \frac{Q_{\text{cond,load}}}{h_2 - h_3}
\end{equation}

응축기 열량($Q_{\text{ref,cond}}$), 증발기 열량($Q_{\text{ref,evap}}$), 압축기 전력($E_{\text{cmp}}$)은 각각 다음과 같이 계산된다:

\begin{equation}
\label{eq:hpb_q_ref_cond}
Q_{\text{ref,cond}} = \dot{m}_{\text{ref}} (h_2 - h_3)
\end{equation}

\begin{equation}
\label{eq:hpb_q_ref_evap}
Q_{\text{ref,evap}} = \dot{m}_{\text{ref}} (h_1 - h_4)
\end{equation}

\begin{equation}
\label{eq:hpb_e_cmp}
E_{\text{cmp}} = \dot{m}_{\text{ref}} (h_2 - h_1)
\end{equation}

\subsubsection{Heat Exchanger Performance}

응축기와 증발기의 열전달은 LMTD(Logarithmic Mean Temperature Difference) 방법으로 계산된다. 응축기의 LMTD는:

\begin{equation}
\label{eq:hpb_lmtd_cond}
\text{LMTD}_{\text{cond}} = \frac{T_2 - T_3}{\ln\left(\frac{T_2 - T_{\text{tank,w}}}{T_3 - T_{\text{tank,w}}}\right)}
\end{equation}

응축기 열전달량은:

\begin{equation}
\label{eq:hpb_q_lmtd_cond}
Q_{\text{LMTD,cond}} = \text{UA}_{\text{cond}} \cdot \text{LMTD}_{\text{cond}}
\end{equation}

증발기의 경우, 실외 공기와 냉매 간의 열교환을 고려하여 LMTD를 계산하고, 이를 통해 필요한 팬 풍량($\dot{V}_{\text{fan,ou}}$)과 실제 열전달량($Q_{\text{LMTD,evap}}$)을 결정한다.

\subsubsection{Outdoor Unit Fan Power}

실외기 팬 전력은 VSD(Variable Speed Drive) 곡선을 기반으로 계산된다:

\begin{equation}
\label{eq:hpb_e_fan}
E_{\text{fan,ou}} = f(\dot{V}_{\text{fan,ou}})
\end{equation}

여기서 $f$는 VSD 곡선 함수이며, 설계 풍량과 설계 전력을 기준으로 비례 관계를 따른다.

\subsubsection{Hot Water Tank Energy Balance}

저탕조의 에너지 밸런스는 다음과 같이 표현된다:

\begin{equation}
\label{eq:hpb_tank_energy}
C_{\text{tank}} \frac{{\rm d}T_{\text{tank,w}}}{{\rm d}t} = Q_{\text{ref,cond}} - Q_{\text{tank,loss}} - Q_{\text{use,loss}}
\end{equation}

탱크 열손실($Q_{\text{tank,loss}}$)과 사용 열손실($Q_{\text{use,loss}}$)은 각각:

\begin{equation}
\label{eq:hpb_q_tank_loss}
Q_{\text{tank,loss}} = \text{UA}_{\text{tank}} (T_{\text{tank,w}} - T_0)
\end{equation}

\begin{equation}
\label{eq:hpb_q_use_loss}
Q_{\text{use,loss}} = c_w \rho_w \dot{V}_{w,\text{sup,tank}} (T_{\text{tank,w}} - T_{\text{sup,w}})
\end{equation}

여기서 $\dot{V}_{w,\text{sup,tank}} = \alpha \dot{V}_{w,\text{serv}}$이며, $\alpha$는 믹싱밸브의 혼합비로 가스 보일러와 동일한 방식으로 계산된다.

\subsubsection{Refrigerant Loop Exergy Analysis}

냉매루프의 엑서지 분석을 위해 각 상태점의 엑서지를 계산한다. 단위 질량당 엑서지($x_i$)는:

\begin{equation}
\label{eq:hpb_exergy_states}
x_i = (h_i - h_0) - T_0(s_i - s_0)
\end{equation}

여기서 $h_0$와 $s_0$는 기준 상태(환경 온도 $T_0$, 환경 압력 $P_0$)에서의 엔탈피와 엔트로피이다. 엑서지 흐름($X_i$)은:

\begin{equation}
\label{eq:hpb_X_i}
X_i = \dot{m}_{\text{ref}} x_i
\end{equation}

컴포넌트별 엑서지 소비는 다음과 같이 계산된다:

\begin{equation}
\label{eq:hpb_xc_evap}
X_{c,\text{evap}} = (X_{a,\text{ou,in}} + X_4) - X_1
\end{equation}

\begin{equation}
\label{eq:hpb_xc_cmp}
X_{c,\text{cmp}} = (X_1 + X_{\text{cmp}}) - X_2
\end{equation}

\begin{equation}
\label{eq:hpb_xc_ref_cond}
X_{c,\text{ref,cond}} = X_2 - X_3 - X_{\text{ref,cond}}
\end{equation}

\begin{equation}
\label{eq:hpb_xc_exp}
X_{c,\text{exp}} = X_3 - X_4
\end{equation}

여기서 $X_{\text{ref,cond}} = (1 - T_0 / T_{\text{tank,w}}) Q_{\text{ref,cond}}$는 응축기에서 저탕조로 전달되는 엑서지이며, $X_{a,\text{ou,in}}$은 실외 공기 입구 엑서지이다.

\subsubsection{Tank and Mixing Valve Exergy Analysis}

저탕조에서의 엑서지 소비는:

\begin{equation}
\label{eq:hpb_xc_tank}
X_{c,\text{tank}} = X_{\text{ref,cond}} + X_{w,\text{sup,tank}} - X_{w,\text{tank}}
\end{equation}

여기서 $X_{w,\text{sup,tank}}$와 $X_{w,\text{tank}}$는 각각 급수와 저탕조 온수의 엑서지이며, 가스 보일러와 동일한 방식으로 계산된다.

믹싱밸브에서의 엑서지 소비는:

\begin{equation}
\label{eq:hpb_xc_mix}
X_{c,\text{mix}} = X_{w,\text{tank}} + X_{w,\text{sup,mix}} - X_{w,\text{serv}}
\end{equation}

\subsubsection{System Exergy Efficiency}

히트펌프 보일러 시스템 전체의 엑서지 효율($\eta_{\text{X,HPB}}$)은 다음과 같이 정의된다:



\begin{equation}
\label{eq:hpb_exergy_efficiency}
\eta_{\text{X,HPB}} = \frac{X_{w,\text{serv}}}{X_{\text{cmp}} + X_{\text{fan,ou}}}
\end{equation}

여기서 $X_{\text{cmp}} = E_{\text{cmp}}$ 및 $X_{\text{fan,ou}} = E_{\text{fan,ou}}$는 각각 압축기와 팬의 엑서지 입력이다. 총 엑서지 소비는:

\begin{equation}
\label{eq:hpb_xc_tot}
X_{c,\text{tot}} = X_{c,\text{evap}} + X_{c,\text{cmp}} + X_{c,\text{ref,cond}} + X_{c,\text{exp}} + X_{c,\text{tank}} + X_{c,\text{mix}}
\end{equation}





본 연구에서는 한국의 급탕 시스템 환경 영향을 정확히 평가하기 위해, 실제 발전 믹스 데이터를 기반으로 1차 에너지 계수(PEF)와 CO$_2$ 배출계수를 산정하였다.

Table~\ref{tab:grid_mix_korea}는 2024년 기준 한국의 전력 발전원별 비중과 각 에너지원의 물리적 PEF 및 CO$_2$ 배출계수를 보여준다. 이를 근거로 한국 전력망의 가중 평균 PEF를 계산하면 다음과 같다.

\begin{equation}
PEF_{grid} = \sum (Share_i \times PEF_i) \approx 2.41
\end{equation}

계산 결과인 \textbf{2.41}은 한국의 건물 에너지 효율 등급 인증 등에서 사용되는 법적 고시값인 \textbf{2.75}보다 약 12\% 낮은 수치이다. 이러한 차이가 발생하는 주된 이유는 다음과 같다:
\begin{enumerate}
    \item \textbf{보수적 산정 (Conservative Estimation):} 법적 고시값은 발전소의 노후화, 송배전 손실 등을 보수적으로 고려하여 설정되는 경향이 있다.
    \item \textbf{최신 효율 미반영:} 최근 도입된 고효율 LNG 복합화력 발전소(PEF $\approx$ 1.8)의 비중이 증가하고 재생에너지 보급이 확대되었으나, 고시값은 이러한 최신 트렌드를 즉각적으로 반영하지 못하는 시차가 존재한다.
\end{enumerate}

본 연구에서는 물리적 실체에 더 근접한 평가를 위해, Table~\ref{tab:grid_mix_korea}에 기반한 계산값(PEF 2.41)을 기준 시나리오(Base case)로 설정하고, 이후 재생에너지(Renewable) 비중을 단계적으로 증가시키는 시나리오 분석을 수행한다.

\begin{table}[htbp]
\centering
\caption{South Korea grid mix energy sources with generation share, primary energy factors, and CO$_2$ emission factors (2024, corrected based on actual generation).}
\label{tab:grid_mix_korea}
\renewcommand{\arraystretch}{1.2}
\small
\begin{tabularx}{\columnwidth}{l X X X}
\toprule
\textbf{Energy Source} & \textbf{Generationshare} & \textbf{Primary energy factor} & \textbf{CO$_2$ Emission factor} \\
 & \textbf{[TWh (\%)]} & \textbf{[--]} & \textbf{[kg CO$_2$/kWh]} \\
\midrule
Coal & 198 (33\%) & 2.6 & 0.90 \\
Gas & 151 (25\%) & 1.8 & 0.40 \\
Nuclear & 189 (31\%) & 3.0 & 0.005 \\
Oil & 7.0 (1\%) & 2.5 & 0.70 \\
Hydro & 9.0 (1.5\%) & 1.0 & 0.01 \\
Wind & 3.4 (0.6\%) & 1.0 & 0.01 \\
Solar & 33.4 (5.5\%) & 1.0 & 0.04 \\
Biomass & 11.8 (2.0\%) & 3.0 & 0.00 \\
\bottomrule
\end{tabularx}
\end{table}

가스 보일러 시스템은 저탕조 없이 직접 급탕 공급 구조로 설계되었다. 이는 현대 가스 보일러의 일반적인 운전 방식인 즉시 가열(on-demand heating) 방식을 반영한 것이다. 시스템은 연소실(combustion chamber)과 믹싱밸브(mixing valve) 두 개의 주요 서브시스템으로 구성된다. 연소실에서는 천연가스를 연소하여 급수를 설정 온도로 가열하고, 믹싱밸브에서는 보일러 출수와 상수도를 혼합하여 서비스 급탕 온도로 공급한다.

공기원 히트펌프 보일러 시스템은 실외기(outdoor unit), 냉매루프(refrigerant loop), 저탕조(hot water tank), 믹싱밸브로 구성된다. 본 연구에서는 냉매 상태를 최적화 알고리즘을 통해 결정하는 방법을 채택하였다. 이는 고정된 운전 조건을 가정하는 기존 연구와 차별화되는 점이다. 히트펌프 사이클의 최적 운전점은 총 전력 사용률($E_{\text{tot}} = E_{\text{cmp}} + E_{\text{fan,ou}}$)를 최소화하는 최적화 문제를 통해 결정된다.

다이어그램에 표시된 변수들과 분석에 사용된 설정값은 Table~\ref{tab:system_parameters}에 정리되어 있다.

\begin{table}[htbp]
\centering
\caption{Assumed parameters for gas boiler (GB) and heat pump boiler (HPB)}
\label{tab:system_parameters}
\renewcommand{\arraystretch}{1.2}
\small
\begin{tabularx}{\columnwidth}{l X c c}
\toprule
\textbf{Symbol [units]} & \textbf{Description} & \textbf{GB} & \textbf{HPB} \\
\midrule
$T_0$ [°C] & Environmental temperature & 0 & 0 \\
$T_{\text{sup,w}}$ [°C] & Temperature of supply water & 15 & 15 \\
$T_{\text{tank,w}}$ [°C] & Temperature of water in hot water tank & -- & 60 \\
$T_{\text{serv,w}}$ [°C] & Temperature of service hot water & 40 & 40 \\
$T_{\text{exh}}$ [°C] & Temperature of exhausted gas from combustion chamber & 70 & -- \\
$T_{\text{comb}}$ [°C] & Combustion chamber setpoint temperature & 60 & -- \\
$T_{\text{a,ou,in}}$ [°C] & Inlet air temperature of outdoor unit & -- & 0 \\
$A_{\text{ext}}$ [m$^2$] & Fan discharge area of outdoor unit & -- & 0.50 \\
$\Delta P_{\text{fan}}$ [Pa] & Pressure difference across outdoor unit & -- & 500 \\
$UA_{\text{tank}}$ [W/K] & Overall heat transfer coefficient of tank wall & -- & 0.45 \\
$\eta_{\text{comb}}$ [--] & Efficiency of combustion chamber & 0.9 & -- \\
$\eta_{\text{NG}}$ [--] & Exergy efficiency of natural gas combustion & 0.93 & -- \\
$\eta_{\text{fan}}$ [--] & Efficiency of fan & -- & 0.8 \\
$\text{UA}_{\text{cond}}$ [W/K] & Overall heat transfer coefficient of condenser & -- & 500 \\
$\text{UA}_{\text{evap}}$ [W/K] & Overall heat transfer coefficient of evaporator & -- & 1500 \\
$V_{\text{disp,cmp}}$ [m$^3$] & Compressor displacement volume per rotation& -- & $5.0 \times 10^{-4}$ \\
$\eta_{\text{cmp,isen}}$ [--] & Isentropic efficiency of compressor & -- & 0.7 \\
\bottomrule
\end{tabularx}
\end{table}

\begin{table}[htbp]
\centering
\caption{Nomenclature}
\label{tab:nomenclature}
\renewcommand{\arraystretch}{1.2}
\begin{tabularx}{\columnwidth}{l X l}
\toprule
\textbf{Symbol} & \textbf{Description} & \textbf{Units} \\
\midrule
$C_{\text{tank}}$ & Thermal capacity of hot water tank & J/K \\
$E_{\text{cmp}}$ & Compressor power consumption & W \\
$E_{\text{fan,ou}}$ & outdoor unit fan power consumption & W \\
$E_{\text{NG}}$ & Natural gas energy input & W \\
$E_{\text{tot}}$ & Total power consumption & W \\
$h_0$ & Enthalpy at reference state & J/kg \\
$h_2$ & Enthalpy at compressor outlet & J/kg \\
$h_3$ & Enthalpy at condenser outlet & J/kg \\
$h_i$ & Enthalpy at state point $i$ & J/kg \\
$\dot{m}_{\text{ref}}$ & Refrigerant mass flow rate & kg/s \\
$Q_{\text{comb,load}}$ & Required heat load based on boiler flow rate & W \\
$Q_{\text{cond,load}}$ & Condenser heat load & W \\
$Q_{\text{exh}}$ & Exhaust gas heat loss & W \\
$Q_{\text{LMTD,cond}}$ & Heat transfer rate in condenser (LMTD method) & W \\
$Q_{\text{LMTD,evap}}$ & Heat transfer rate in evaporator (LMTD method) & W \\
$Q_{\text{ref,cond}}$ & Refrigerant condensation heat transfer rate & W \\
$Q_{\text{ref,evap}}$ & Refrigerant evaporation heat transfer rate & W \\
$Q_{\text{tank,loss}}$ & Hot water tank heat loss & W \\
$Q_{\text{use,loss}}$ & Heat loss due to hot water usage & W \\
$Q_{w,\text{comb,out}}$ & Boiler outlet water heat rate & W \\
$Q_{w,\text{serv}}$ & Service hot water heat rate & W \\
$Q_{w,\text{sup}}$ & Supply water heat rate & W \\
$Q_{w,\text{sup,mix}}$ & Supply water heat rate to mixing valve & W \\
$T_0$ & Environmental temperature (reference temperature) & K \\
$T_{\text{cond}}$ & Condensation temperature & K \\
$T_{\text{evap}}$ & Evaporation temperature & K \\
$T_{\text{exh}}$ & Exhaust gas temperature & K \\
$T_{\text{serv,w}}$ & Service hot water temperature & K \\
$T_{\text{sup,w}}$ & Supply water temperature & K \\
$T_{\text{tank,w}}$ & Hot water tank temperature & K \\
$T_{\text{w,comb}}$ & Boiler outlet water temperature & K \\
$\text{UA}_{\text{cond}}$ & Overall heat transfer coefficient of condenser & W/K \\
$\text{UA}_{\text{evap}}$ & Overall heat transfer coefficient of evaporator & W/K \\
$\dot{V}_{w,\text{sup,comb}}$ & Supply water volumetric flow rate to combustion chamber & m$^3$/s \\
$X_{\text{c}}$ & Exergy consumption rate & W \\
$X_i$ & Exergy flow rate at state point $i$ & W \\
$X_{\text{NG}}$ & Natural gas exergy input & W \\
$X_{w,\text{serv}}$ & Service hot water exergy & W \\
$x_i$ & Specific exergy at state point $i$ & J/kg \\
\vspace{0.1cm} \\
\multicolumn{3}{@{}l}{\textit{Greek Symbols}} \\
$\alpha$ & Mixing ratio & -- \\
$\Delta T_{\text{ref,cond}}$ & Temperature difference between refrigerant and hot water tank & K \\
$\Delta T_{\text{ref,evap}}$ & Temperature difference between refrigerant and outdoor air & K \\
$\epsilon$ & Tolerance (allowable error) & -- \\
$\eta_{\text{comb}}$ & Combustion efficiency & -- \\
$\eta_{\text{X,GB}}$ & Exergy efficiency of gas boiler system & -- \\
$\eta_{\text{X,HPB}}$ & Exergy efficiency of heat pump boiler system & -- \\
\bottomrule
\end{tabularx}
\end{table}


\end{document}



% REVIEW: The observation that Adding more sensors can *degrade* performance if poorly placed is a counter-intuitive and strong finding. It strongly validates the necessity of the proposed placement algorithm. The conclusion calling for "quality over quantity" in data collection is compelling.

\begin{figure*}[htbp]
  \centering
  \begin{tabular}{cc}
    \fbox{\parbox{0.9\columnwidth}{\centering \vspace{2cm} PLACEHOLDER (a) \vspace{2cm}}} &
    \fbox{\parbox{0.9\columnwidth}{\centering \vspace{2cm} PLACEHOLDER (b) \vspace{2cm}}} \\
    (a) Temperature & (b) Humidity \\
    \fbox{\parbox{0.9\columnwidth}{\centering \vspace{2cm} PLACEHOLDER (c) \vspace{2cm}}} &
    \fbox{\parbox{0.9\columnwidth}{\centering \vspace{2cm} PLACEHOLDER (d) \vspace{2cm}}} \\
    (c) CO2 & (d) NH3 \\
  \end{tabular}
  \caption{물리 센서 개수 및 배치에 따른 가상 센서 예측 RMSE 평균 (Red: 최적 배치, Gray: 그 외). (a) 온도, (b) 습도, (c) 이산화탄소, (d) 암모니아.}
  \label{fig:perf_trends}
\end{figure*}

% REVIEW: The observation that non-optimal placement of additional sensors can degrade average performance is a key contribution. It contradicts the common assumption that "more data is better" and highlights the importance of the proposed optimization framework. Ensure this point is highlighted in the conclusion as well.



%% Use \section commands to start a section
\section{Conclusions}


%% Use \subsection commands to start a subsection.
\subsection{aaaa}



% \end{linenumbers}

\section*{Acknowledgments}
This work was supported by the National Research Foundation of Korea (NRF) grant (No. RS-2023-00277318 and RS-2025-00512551) funded by the Korean government (Ministry of Science and ICT).


상호 참조하기 -> Fig. \ref{fig:fig1}
%% The Appendices part is started with the command \appendix;
%% appendix sections are then done as normal sections
\appendix
\renewcommand{\figurename}{Fig} % "Figure"를 "Fig"로 변경
\renewcommand{\thefigure}{\Alph{section}\arabic{figure}} % "Appendix B" 대신 "B"만 가져옴
\renewcommand{\thetable}{\Alph{section}\arabic{table}} % 표 번호도 동일하게 적용

%% Appnedix A
\section{Appendix A. aaaa}
\label{appendA}
\setcounter{figure}{0}
\setcounter{table}{0}


Appendix text.


%% Appendix B 
\section{Appendix B. bbbbb}
\label{appendB}
\setcounter{figure}{0}
\setcounter{table}{0}

Appendix B text 

\begin{figure}[htbp]%% placement specifier
%% Use \includegraphics command to insert graphic files. Place graphics files in 
%% working directory.
\centering%% For centre alignment of image.
%% Use \caption command for figure caption and label.
\caption{{컬럼 너비 피겨 (옵션 \texttt{width=\textbackslash columnwidth} 사용)}}
\label{fig:appendix-fig1}
%% https://en.wikibooks.org/wiki/LaTeX/Importing_Graphics#Importing_external_graphics
\end{figure}


%% Supplementary material 섹션
\section*{Supplementary material}
Interactive plot for time-series pattern exploration: An interactive HTML extension of Fig. \ref{fig:humidity_same_type}(a) is provided with two key features

%% Code and data availability 섹션
\section*{Code and data availability}



%% If you have bib database file and want bibtex to generate the
%% bibitems, please use
%%

% \bibliographystyle{elsarticle-num} 
\bibliographystyle{elsarticle-harv} 
\bibliography{references}

%%  \bibliography{<your bibdatabase>}

%% else use the following coding to input the bibitems directly in the
%% TeX file.

%% Refer following link for more details about bibliography and citations.
%% https://en.wikibooks.org/wiki/LaTeX/Bibliography_Management



%% Examples 섹션은 필요시 주석 해제하여 사용하세요
% \section{Examples}
% \subsection{Figures}
% \begin{figure}[htbp]
%   \centering
%   \includegraphics[width=\linewidth]{figure_name} % 확장자(.eps, .pdf) 생략 가능
%   \caption{Description of the figure.}
%   \label{fig:my_figure}
% \end{figure}



\end{document}
